\documentclass[10pt,a4paper]{amsart}
\usepackage[margin=19mm]{geometry}
\usepackage{amsmath,amssymb,mathtools}
\usepackage[scr=rsfs,cal=euler]{mathalfa}
\usepackage{xfrac}
\usepackage[unicode,hidelinks,bookmarks=false]{hyperref}
\newcommand{\ie}{i.e.\ }
\newcommand{\cf}{cf.\ }
\newcommand{\resp}{respectively\ }
\newcommand{\Subsection}{\subsection}
\providecommand{\nobreakdash}{\nobreak\hbox{-}}
\setlength{\emergencystretch}{2em}
\multlinegap0pt
\usepackage{bm}
\usepackage{bbm}
\usepackage{dsfont}
\usepackage{xspace}
\usepackage{cancel}
\usepackage{mathtools}
\usepackage{amsthm}
\makeatletter
\@ifundefined{defn}{\newtheorem{defn}{Defn}}{}
\@ifundefined{assumption}{\newtheorem{assumption}[defn]{Assumption}}{}
\@ifundefined{lemma}{\newtheorem{lemma}[defn]{Lemma}}{}
\makeatother
\title[Nonsmooth bifurcations in families of one-dimensional piecew]{Nonsmooth bifurcations in families of one-dimensional piecewise-linear quasiperiodically forced maps}
\author{Rafael Martinez-Vergara, Joan Carles Tatjer}
\date{Source: arXiv:2512.04234v2 (CC BY 4.0)}
\begin{document}
\maketitle

\noindent\emph{Source and licence.} Nonsmooth bifurcations in families of one-dimensional piecewise-linear quasiperiodically forced maps. Version arXiv:2512.04234v2 (CC BY 4.0), distributed under the Creative Commons Attribution 4.0 International licence, as recorded in the arXiv licence metadata for this exact version and shown on the abstract page https://arxiv.org/abs/2512.04234v2. Full rights notice: licenses/arxiv-2512.04234-CC-BY-4.0.txt.\par\smallskip

\noindent\textit{Editorial scope.} Counted unit: the Lemma whose source label is \texttt{zeros\_lambda} in the article (source0.tex lines 967--973, source label \texttt{zeros\_lambda}), delivered together with the article's own complete original proof of it (source0.tex lines 974--1017, one \texttt{proof} environment of 44 lines). No mathematics is written, repaired or re-derived: statement and proof are the article's own text line for line. Required context retained uncounted: dyn\_sys (lines 152--164); dyn\_sys\_2 (lines 168--180); dyn\_sys\_3 (lines 184--195); dyn\_sys\_4 (lines 199--210); assumptions\_a (lines 288--290); expression\_lambda\_4 (lines 766--772). Every definition, display and label of those lines is retained unchanged. Omitted from this package and not redistributed: 1--151, 165--167, 181--183, 196--198, 211--287, 291--765, 773--966, 1018--1470. Nothing inside a retained excerpt is omitted or altered: every retained body is delivered exactly as the article's lines are written, so no omission receipt of any kind accompanies this package. Citations: the retained excerpts carry no citation command at all, so no bibliography entry of the article is transcribed and no reference is redistributed. Reference adaptation: none was needed and none was made; every reference inside the delivered unit resolves inside it, so no unresolved reference or citation is produced. Printed slips kept literal: no printed word, symbol, hypothesis or number of the counted statement or of its proof was altered. Layout and class: the article's own class is \texttt{article} with packages including graphicx, amsmath, amssymb, babel, amsthm, mathtools, xcolor, geometry, setspace, hyperref, comment, enumitem, subcaption, caption; the wrapper is the lane's pinned \texttt{amsart} editorial wrapper at 10pt on A4, which restates the theorem-like environments the retained text needs and adds the article's own macro definitions that the retained text needs (none), with extra wrapper packages (bm, bbm, dsfont, xspace, cancel, mathtools). The delivered numbering is the unit's own. Assets: no retained span contains a figure environment, a graphics-inclusion command, a PDF-page inclusion, a table or a TikZ picture; no image file is copied into this package, the package declares no asset dependency and no image file is redistributed with it. Licence: version arXiv:2512.04234v2 of this article is distributed under the Creative Commons Attribution 4.0 International licence, as recorded in the arXiv licence metadata for this exact version and shown on the abstract page https://arxiv.org/abs/2512.04234v2. A full rights notice is delivered at licenses/arxiv-2512.04234-CC-BY-4.0.txt. Characters: every retained excerpt is pure ASCII, so the delivered engine, XeLaTeX with its default Latin Modern text font, reports no missing character.
    \begin{align}\label{dyn_sys}
        \begin{cases}
            \Bar{x} = h_{1}(x) - b g(\theta),\\
            \Bar{\theta} = \theta + \omega \mod2\pi,
        \end{cases}
        &\hspace{1cm}
        h_{1}(x)= 
        \begin{cases}
            -\frac{\pi}{2} &\text{if }x\in(-\infty,-\frac{\pi}{2a}),\\
            ax & \text{if }x\in[-\frac{\pi}{2a},\frac{\pi}{2a}],\\
            \frac{\pi}{2} &\text{if }x\in(\frac{\pi}{2a},\infty),
        \end{cases}
    \end{align}

    \begin{align}\label{dyn_sys_2}
        \begin{cases}
            \Bar{x} = h_{2}(x) + b g(\theta),\\
            \Bar{\theta} = \theta + \omega \mod2\pi,
        \end{cases}
        &\hspace{1cm}
        h_{2}(x)= 
        \begin{cases}
            a(x+\delta) &\text{if }x\in(-\infty,-\delta),\\
            0 & \text{if }x\in[-\delta,\delta],\\
            a(x-\delta) &\text{if }x\in(\delta,\infty),
        \end{cases}
    \end{align}

    \begin{align}\label{dyn_sys_3}
        \begin{cases}
            \Bar{x} = h_{3}(x) + b g(\theta),\\
            \Bar{\theta} = \theta + \omega \mod2\pi,
        \end{cases}
        &\hspace{1cm}
        h_{3}(x)= 
        \begin{cases}
            ax & \text{if }x>-\frac{1}{a},\\
            -1 &\text{if }x\leq-\frac{1}{a},
        \end{cases}
    \end{align}

    \begin{align}\label{dyn_sys_4}
        \begin{cases}
            \Bar{x} = h_{4}(x) - b g(\theta),\\
            \Bar{\theta} = \theta + \omega \mod2\pi,
        \end{cases}
        &\hspace{1cm}
        h_{4}(x)= 
        \begin{cases}
            ax & \text{if }x>\frac{1}{a},\\
            1 &\text{if }x\leq \frac{1}{a},
        \end{cases}
    \end{align}

\begin{assumption}\label{assumptions_a}
    For the systems \eqref{dyn_sys}, \eqref{dyn_sys_2} and \eqref{dyn_sys_3} we assume $a>1$ and for the \eqref{dyn_sys_4} we assume $a<-1$.
\end{assumption}

    \begin{align}\label{expression_lambda_4}
        \lambda_{n+1}(\theta) =
        \begin{cases}
            1 - a\mu_4(\theta-\omega) &\text{ if }a[\mu_4(\theta-2\omega)+\lambda_n(\theta-2\omega)] - b g(\theta-2\omega)\leq\frac{1}{a},\\
            a^2\lambda_n(\theta-2\omega) &\text{ if }a[\mu_4(\theta-2\omega)+\lambda_n(\theta-2\omega)] - b g(\theta-2\omega)>\frac{1}{a}.
        \end{cases}
    \end{align}

\begin{lemma}\label{zeros_lambda}
    Let $a$ satisfy Assumption \eqref{assumptions_a} and $b = b^*(a)$. The curve $\lambda_n$ has a finite set of of zeros. Concretely, let $\theta_0$ be the only zero of $\lambda_0$. Then we have the following:
    \begin{itemize}
        \item For the systems \eqref{dyn_sys}, \eqref{dyn_sys_2} and \eqref{dyn_sys_3}, $\lambda_n$ has exactly $n+1$ zeros $\theta_0,\theta_0+\omega,\dots,\theta_0+n\omega$.
        \item For the system \eqref{dyn_sys_4}, $\lambda_n$ has exactly $n+1$ zeros $\theta_0,\theta_0+2\omega,\dots,\theta_0+2n\omega$.
    \end{itemize}
\end{lemma}

\begin{proof}
    Consider the system \eqref{dyn_sys_4}.
    Given $n\geq 1$, we see that
    \begin{align*}
        F_{4}^{2(n-1)}(\mu(\theta_0-2n\omega),\theta_0-2n\omega) &= (\mu(\theta_0-2\omega),\theta_0-2\omega),\\
        F_{4}^{2(n-1)}(\varphi_0(\theta_0-2n\omega),\theta_0-2n\omega) &= (\varphi_{n-1}(\theta_0-2\omega),\theta_0-2\omega).
    \end{align*}
    Since $\mu(\theta_0-\omega)=\frac{1}{a}$ and $\mu(\theta_0-2\omega)\leq \varphi_{n-1}(\theta_0-2\omega)$,
    \begin{align*}
        h_4(\varphi_{n-1}(\theta_0-2\omega))-b g(\theta)\leq h_4(\mu(\theta_0-2\omega))-b g(\theta) = \mu(\theta_0-\omega) \leq \frac{1}{a}.
    \end{align*}
    Therefore, we obtain that
    \begin{align*}
        F_{4}^2(\mu(\theta_0-2\omega),\theta_0-2\omega) = F_{4}^2(\varphi_{n-1}(\theta_0-2\omega),\theta_0-2\omega).
    \end{align*}
    As a consequence, we obtain that $\lambda_n(\theta_0)=\varphi_n(\theta_0)-\mu(\theta_0)=0$. If $\theta$ is a zero of $\lambda_n$, from \eqref{expression_lambda_4} we have that
    \begin{align*}
        \lambda_{n+1}(\theta+2\omega) = 
        \begin{cases}
            1 - a\mu(\theta+\omega) &\text{ if }a\mu(\theta) - b g(\theta)\leq\frac{1}{a},\\
            a^2\lambda_n(\theta) &\text{ if }a\mu(\theta) - b g(\theta)>\frac{1}{a}.
        \end{cases}
    \end{align*}
    Since $a\mu(\theta) - b g(\theta) = \mu(\theta+\omega)\geq\frac{1}{a}$, $\lambda_{n+1}(\theta+2\omega)=a^2\lambda_n(\theta)$. Then $\theta+2\omega$ is a zero of $\lambda_{n+1}$. 
    If we consider the case where $\theta$ is the minimum of $\mu$, then $\theta=\theta_0-\omega$ so
    \begin{align*}
        \lambda_{n+1}(\theta+\omega) = 
        \begin{cases}
            1 - a\mu(\theta) &\text{ if }a\mu(\theta-\omega) - b g(\theta-\omega)\leq\frac{1}{a},\\
            a^2\lambda_n(\theta-\omega) &\text{ if }a\mu(\theta-\omega) - b g(\theta-\omega)>\frac{1}{a}.
        \end{cases}
    \end{align*}
    Since $a\mu(\theta-\omega) - b g(\theta-\omega) = \mu(\theta) = \frac{1}{a}$, $\lambda_{n+1}(\theta+\omega) = 1 - a\mu(\theta) = 0$.
    
    The uniqueness follows from induction. For $n=0$ this is true, so we assume that $\lambda_{n+1}$ has an extra zero $\theta_1$. From \eqref{expression_lambda_4} there are two cases,
    \begin{enumerate}
        \item $1 - a\mu(\theta+\omega)=0$.
        We have that $\mu(\theta_1-\omega)=\frac{1}{a}$ so $\theta_1=\theta_0+2\omega$.
        
        \item $a^2\lambda_n(\theta)=0$.
        In this case $\lambda_n$ has an extra zero, but this is contradictory with the induction hypothesis which is that $\lambda_{n}$ has no extra zero.
    \end{enumerate}
    A similar argument is valid for the systems \eqref{dyn_sys}, \eqref{dyn_sys_2} and \eqref{dyn_sys_3}. For these systems you have to consider the one-iteration map instead of the two-iteration map.
\end{proof}

\end{document}
